\chapter{Charakterystyka projektowanego stanowiska laboratoryjnego}
\label{chap: Charakterystyka projektowanego stanowiska laboratoryjnego}

W~rozdziale scharakteryzowano projektowane stanowisko laboratoryjne
z~jednofazowym falownikiem napięcia. Określono wymagania funkcjonalne,
parametry elektryczne układu, wymagania dotyczące bezpieczeństwa
użytkowania oraz wymagania dydaktyczne, a~następnie przedstawiono
architekturę stanowiska, wynikającą z~tych wymagań. Sformułowane tu
założenia stanowią punkt wyjścia do doboru elementów obwodu mocy, opisanego
w~rozdziale~\ref{chap: Dobór elementów obwodu mocy falownika}.

% ---------------------------------------------------------------------
\section{Wymagania funkcjonalne}
\label{sec: Wymagania funkcjonalne}

Projektowane stanowisko ma umożliwić praktyczne zapoznanie się z~pracą
jednofazowego falownika napięcia w~układzie mostka~H oraz doświadczalną
weryfikację zjawisk omówionych w~rozdziale~\ref{chap: Falowniki jednofazowe - podstawy teoretyczne}.
Stanowi ono uzupełnienie zajęć symulacyjnych z~falowników
jednofazowych~\cite{instrukcja}, w~trakcie których badane są modele
falowników sterowanych metodą przebiegu prostokątnego, metodą modulacji
szerokości pojedynczego impulsu w~półokresie oraz metodami bipolarnej
i~unipolarnej modulacji SPWM. Stanowisko powinno pozwalać na
odtworzenie tych samych przebiegów w~rzeczywistym układzie, a~także na
obserwację zjawisk pomijanych w~modelach idealnych, takich jak procesy
łączeniowe tranzystorów, wpływ czasu martwego, straty mocy czy zakłócenia
związane z~szybkimi zmianami napięcia.

W~wymaganiach funkcjonalnych określono, \emph{jakie funkcje} ma realizować
stanowisko, bez wskazywania konkretnych wartości liczbowych. Parametry
elektryczne zestawiono w~podrozdziale~\ref{sec: Parametry elektryczne projektowanego układu},
a~szczegółowe wymagania dotyczące bezpieczeństwa i~dydaktyki --
w~podrozdziałach~\ref{sec: Wymagania dotyczące bezpieczeństwa użytkowania}
i~\ref{sec: Wymagania dydaktyczne}. Przyjęto następujące wymagania
funkcjonalne.

\begin{enumerate}[label=\textbf{F\arabic*.}, leftmargin=*, itemsep=4pt]

  \item \textbf{Przekształcanie energii prądu stałego na energię prądu
    przemiennego.} Stanowisko ma pracować jako falownik napięcia w~układzie
    pełnego mostka (dwie gałęzie łącznikowe, tranzystory $Q_1$--$Q_4$).
    Obwód pośredniczący należy zasilać z~zewnętrznego zasilacza
    laboratoryjnego o~regulowanym napięciu, a~nie bezpośrednio z~sieci
    elektroenergetycznej. Pozwala to rozpoczynać pracę od niskiego napięcia
    $U_\mathrm{DC}$ i~zwiększać je stopniowo oraz ograniczać prąd pobierany
    przez układ nastawą zasilacza. Ponieważ zasilacz laboratoryjny zwykle
    nie może przyjmować energii, energię zwracaną z~obciążenia
    o~charakterze indukcyjnym musi przejmować kondensator obwodu
    pośredniczącego $C_\mathrm{DC}$.

  \item \textbf{Niezależne sterowanie każdym z~czterech tranzystorów.}
    Każdy tranzystor mostka ma mieć własny, galwanicznie odizolowany kanał
    sterowania bramkowego, zasilany z~osobnego izolowanego źródła. Takie
    rozwiązanie, w~odróżnieniu od zasilania typu \textit{bootstrap}, nie
    nakłada ograniczeń na współczynnik wypełnienia ani na czas, przez jaki
    tranzystor górny może pozostawać załączony. Umożliwia to realizację
    dowolnego schematu łączeń, w~tym sterowania przebiegiem prostokątnym
    o~częstotliwości $f_1$, w~którym każdy tranzystor przewodzi przez
    połowę okresu napięcia wyjściowego, a~także pracę statyczną, np.~przy
    sprawdzaniu poszczególnych tranzystorów.

  \item \textbf{Realizacja różnych metod sterowania.} Stanowisko ma
    umożliwiać sterowanie falownikiem co najmniej metodami opisanymi
    w~podrozdziale~\ref{sec: Metody modulacji PWM}: przebiegiem
    prostokątnym, modulacją szerokości pojedynczego impulsu (przez
    przesunięcie fazowe sygnałów sterujących gałęzi) oraz bipolarną
    i~unipolarną modulacją SPWM. Zmiana metody sterowania nie powinna
    wymagać zmian w~obwodzie mocy, a~jedynie zmiany sygnałów sterujących.

  \item \textbf{Regulacja parametrów sterowania.} Użytkownik ma mieć
    możliwość zmiany częstotliwości podstawowej $f_1$, współczynnika
    modulacji amplitudowej $m_\mathrm{a}$, częstotliwości łączeń
    $f_\mathrm{s}$ oraz czasu martwego $t_\mathrm{d}$. Regulacja
    $m_\mathrm{a}$ powinna obejmować również zakres nadmodulacji, a~regulacja
    $t_\mathrm{d}$ -- wartości wyraźnie większe od minimalnych, tak aby
    można było obserwować zniekształcenia napięcia wyjściowego wywołane
    czasem martwym.

  \item \textbf{Oddzielenie części sterującej od obwodu mocy.} Sygnały
    sterujące tranzystorami mają być doprowadzane do stanowiska jako
    sygnały logiczne z~zewnętrznego układu sterującego (mikrokontrolera lub
    generatora sygnałów). Wejścia sterujące oraz obwody pomiarowe
    i~zabezpieczające po stronie układu sterującego mają być galwanicznie
    odizolowane od obwodu mocy. Dzięki temu układ sterujący, komputer
    i~przyrządy pomiarowe użytkownika nie są narażone na potencjały
    występujące w~mostku.

  \item \textbf{Zasilanie różnych odbiorników.} Do wyjścia falownika ma
    być możliwe dołączenie obciążenia rezystancyjnego,
    rezystancyjno-indukcyjnego oraz filtru wyjściowego LC. Pozwala to
    porównać przebiegi prądu wyjściowego $i_\mathrm{o}$ i~napięcia
    wyjściowego $u_\mathrm{o}$ przed filtrem i~za nim oraz zbadać wpływ
    indukcyjności $L_\mathrm{f}$ na tętnienia prądu.

  \item \textbf{Dostęp do sygnałów pomiarowych.} Stanowisko ma zapewniać
    łatwy dostęp do punktów pomiarowych, umożliwiających rejestrację
    oscyloskopem co najmniej: napięć gałęzi $u_\mathrm{aN}$
    i~$u_\mathrm{bN}$, napięcia wyjściowego $u_\mathrm{o}$, prądu
    wyjściowego $i_\mathrm{o}$, napięcia obwodu pośredniczącego
    $U_\mathrm{DC}$ oraz sygnałów sterujących tranzystorami. Prąd
    wyjściowy ma być mierzony w~układzie wbudowanym, dającym sygnał
    napięciowy proporcjonalny do prądu, co pozwala obserwować go bez
    sondy prądowej.

  \item \textbf{Możliwość pomiaru sprawności.} Zaciski wejściowe
    i~wyjściowe falownika mają umożliwiać dołączenie przyrządów do pomiaru
    mocy pobieranej z~zasilacza i~mocy oddawanej do obciążenia, tak aby
    można było wyznaczyć sprawność przekształtnika w~funkcji obciążenia.

  \item \textbf{Zabezpieczenia i~sygnalizacja stanu.} Stanowisko ma
    samoczynnie blokować pracę tranzystorów po przekroczeniu
    dopuszczalnego prądu, niezależnie od oprogramowania układu
    sterującego, oraz informować układ sterujący o~zadziałaniu
    zabezpieczenia. Tranzystory nie mogą być załączane, gdy napięcie
    zasilania sterowników bramkowych jest zbyt niskie. Po wyłączeniu
    stanowiska musi być możliwe szybkie rozładowanie kondensatora obwodu
    pośredniczącego przez rezystory rozładowcze, załączane przyciskiem,
    a~stan pracy układu (obecność napięcia, zadziałanie zabezpieczeń) ma
    być sygnalizowany optycznie. Szczegółowe wymagania
    przedstawiono w~podrozdziale~\ref{sec: Wymagania dotyczące bezpieczeństwa użytkowania}.

  \item \textbf{Łatwa wymiana łączników.} Tranzystory mocy mają być
    montowane w~sposób umożliwiający ich wymianę bez lutowania. Pozwala to
    szybko wymienić uszkodzony element, a~także porównywać pracę
    falownika z~tranzystorami różnych typów.

\end{enumerate}

Zestawienie wymagań funkcjonalnych wraz ze wskazaniem miejsc pracy,
w~których opisano sposób ich realizacji i~weryfikacji, przedstawiono
w~tabeli~\ref{tab:wymagania-funkcjonalne}.

\begin{table}[htb]
  \centering
  \caption{Zestawienie wymagań funkcjonalnych projektowanego stanowiska}
  \label{tab:wymagania-funkcjonalne}
  \small
  \begin{tabular}{@{}C{10mm} L{80mm} L{55mm}@{}}
    \toprule
    \textbf{Nr} & \textbf{Wymaganie} & \textbf{Realizacja / weryfikacja} \\
    \midrule
    F1  & Przekształcanie DC/AC w~układzie mostka~H, zasilanie z~zasilacza
          laboratoryjnego
        & podrozdz.~\ref{sec: Architektura projektowanego stanowiska},
          \ref{sec: Dobór kondensatorów magistrali DC} \\
    F2  & Niezależne, izolowane sterowanie każdego z~tranzystorów $Q_1$--$Q_4$
        & podrozdz.~\ref{sec: Dobór sterowników bramkowych} \\
    F3  & Sterowanie przebiegiem prostokątnym, pojedynczym impulsem oraz
          SPWM bipolarną i~unipolarną
        & podrozdz.~\ref{sec: Analiza przebiegów napięć i prądów},
          \ref{sec: Przebiegi napięć i prądów} \\
    F4  & Regulacja $f_1$, $m_\mathrm{a}$, $f_\mathrm{s}$ i~$t_\mathrm{d}$
        & podrozdz.~\ref{sec: Możliwe ćwiczenia laboratoryjne} \\
    F5  & Izolacja galwaniczna części sterującej od obwodu mocy
        & podrozdz.~\ref{sec: Architektura projektowanego stanowiska},
          \ref{sec: Projekt obwodu drukowanego} \\
    F6  & Obciążenie R, RL oraz filtr wyjściowy LC
        & podrozdz.~\ref{sec: Przebiegi napięć i prądów} \\
    F7  & Punkty pomiarowe napięć, prądu i~sygnałów sterujących
        & podrozdz.~\ref{sec: Projekt obwodu drukowanego},
          \ref{sec: Stanowisko pomiarowe} \\
    F8  & Pomiar sprawności przekształtnika
        & podrozdz.~\ref{sec: Badanie sprawności przekształtnika} \\
    F9  & Zabezpieczenie nadprądowe i~podnapięciowe, rozładowanie
          $C_\mathrm{DC}$, sygnalizacja stanu
        & podrozdz.~\ref{sec: Wymagania dotyczące bezpieczeństwa użytkowania},
          \ref{sec: Ocena działania zabezpieczeń} \\
    F10 & Wymiana tranzystorów bez lutowania
        & podrozdz.~\ref{sec: Projekt obwodu drukowanego} \\
    \bottomrule
  \end{tabular}
  \zrodlo{opracowanie własne}
\end{table}


\section{Parametry elektryczne projektowanego układu}
\label{sec: Parametry elektryczne projektowanego układu}

\section{Wymagania dotyczące bezpieczeństwa użytkowania}
\label{sec: Wymagania dotyczące bezpieczeństwa użytkowania}

\section{Wymagania dydaktyczne}
\label{sec: Wymagania dydaktyczne}

\section{Architektura projektowanego stanowiska}
\label{sec: Architektura projektowanego stanowiska}
